\documentclass[10pt,a4paper]{amsart}
\usepackage[margin=19mm]{geometry}
\usepackage{amsmath,amssymb,mathtools}
\usepackage[scr=rsfs,cal=euler]{mathalfa}
\usepackage{xfrac}
\usepackage[unicode,hidelinks,bookmarks=false]{hyperref}
\newcommand{\ie}{i.e.\ }
\newcommand{\cf}{cf.\ }
\newcommand{\resp}{respectively\ }
\newcommand{\Subsection}{\subsection}
\providecommand{\nobreakdash}{\nobreak\hbox{-}}
\setlength{\emergencystretch}{2em}
\multlinegap0pt
\usepackage[shortlabels]{enumitem}
\usepackage{amssymb}
\usepackage{dsfont}
\newtheorem{thm}{Theorem}
\title{Study of weaving frames in Krein spaces}
\author{Avinash Bhardwaj, Animesh Bhandari$^*$}
\date{Source: arXiv:2502.20716v1 (CC BY 4.0)}
\begin{document}
\maketitle

\noindent\emph{Source and licence.} Study of weaving frames in Krein spaces. Version arXiv:2502.20716v1 (CC BY 4.0), distributed by arXiv under the Creative Commons Attribution 4.0 International licence, http://creativecommons.org/licenses/by/4.0/, as recorded in the arXiv licence metadata for this exact version and shown on the abstract page https://arxiv.org/abs/2502.20716v1. Full licence notice: \texttt{licenses/arxiv-2502.20716-CC-BY-4.0.txt}.\par\smallskip

\noindent\textit{Editorial scope.} Counted unit: the article's thm thm (source0.tex lines 469--479). The article's complete original proof of it is delivered with it (source0.tex lines 481--522, one \texttt{proof} environment of 42 lines). No mathematics is written, repaired or re-derived: the statement and the proof are the article's own lines, in the article's own notation, and no retained line is substituted or re-flowed.  No further source-local context is needed: every cross-reference of every retained line resolves inside the two delivered spans.  Nothing was omitted from any retained span, so the package carries no omission record.  No retained line contains a citation, so the wrapper carries no bibliography and no bibliography entry.  No retained line contains a figure environment, an image-inclusion command or drawing code, so no image is delivered and the package declares no asset dependency.  Editorial formatting only, in addition: an AMS \texttt{amsart} preamble that restates the article's own declarations and macros used by the retained lines, namely the environments \texttt{thm} on separate counters, together with the standard packages those lines need.  The retained environments print in this document as Theorem 1 in order of appearance, whereas the article prints the same items with the numbering of its own full document.  The complete native source of the article as distributed in the arXiv e-print for this exact version is bundled unchanged as \texttt{sources/174004/source0.tex}.
\begin{thm}
	
	Let $\{z_n\}_{n=1}^\infty$ and $\{z'_n\}_{n=1}^\infty$ be two frames for Krein space $\mathcal{K},[., .]$, then the following are equivalent:
	
	\begin{enumerate}
	
	\item  The frames $\{z_n\}_{n=1}^\infty$ and $\{z'_n\}_{n=1}^\infty$ are weaving frames for $\mathcal{K}$.
	
\item	 For every $\sigma\subset \mathbb{N}$, let $Z_\sigma = \overline{span}\{z_n\}_{n \in\sigma}$,  $Z_{\sigma^{c}}= \overline{span}\{z'_n\}_{n \in\sigma^{c}}$ and let $P$ be the orthogonal projection of $\mathcal{K}$ onto $Z^\perp_{\sigma}$. Then $P|Z_{\sigma^{c}}$ is the orthogonal projection onto $Z^\perp_\sigma$ and $\{Pz'_n\}_ {n \in \sigma^c}$ is a frame for $Z^\perp_{\sigma}$.
		\end{enumerate}
\end{thm}

\begin{proof}
	
 \noindent  (\underline{1 $\implies$ 2})\\
 First we prove that for every $\sigma \subset \mathbb N$, $P|Z_{\sigma^c}$ be the orthogonal projection  onto $Z^\perp_{\sigma^c}$.
 
	Let us define $ Z_{\sigma} = \overline{span} \{z_n\}_{n \in \sigma}$, 
	$ Z_{\sigma^c} = \overline{span} \{z'_n\}_{n \in \sigma^c}$ and P be the orthogonal projection of $\mathcal{K}$ onto $Z_{\sigma}^\perp$.
	Assumed condition implies that for every $\sigma \subset \mathbb N$, $(\{z_n\}_{n \in \sigma} \cup \{z'_n\}_{n \in \sigma^c})$ is a  frame for $\mathcal{K}$.
	Since $\{z'_n\}_{n \in \sigma^c}$ spans $Z_{\sigma}^\perp$ then every $k \in Z_{\sigma}^\perp$ can be written as linear combination of $\{z'_n\}_{n \in \sigma^c}$, therefore  $P|Z_{\sigma^c}$ be the orthogonal projection  onto $Z^\perp_{\sigma}$.
	
%	Now, to show that $\{Pz'_n\}_ {n \in \sigma^c}$ is a frame for $Z_{\sigma^{c}}$.

	Since $(\{z_n\}_{n \in \sigma} \cup \{z'_n\}_{n \in \sigma^c})$ is a  frame for $\mathcal{K}$, then for every  $k \in \mathcal{K}$, there are universal constants $\alpha \leq \beta$ we have,
	$$\alpha \|k\|_{J}^2 \leq \sum\limits_{n \in \sigma} |[k, z_n]|^2 + \sum\limits_{n \in \sigma^c} |[k, z'_n]|^2 \leq \beta \|k\|_{J}^2.$$
Hence	for every $k \in Z_{\sigma}^\perp$ we have,
	$$\alpha \|k\|_{J}^2 \leq  \sum\limits_{n \in \sigma^c} |[k, z'_n]|^2 \leq \beta \|k\|_{J}^2.$$
	Since $P(z'_n) = z'_n$, then $\{Pz'_n\}_{n \in \sigma^c}$ is a frame for $Z_{\sigma}^\perp.$\\
	
	\noindent  (\underline{2 $\implies$ 1})\\
	Since $\{z_n\}_{n=1}^\infty$ is a frame  for $\mathcal{K}$ and $ Z_\sigma = \overline{span}(\{z_n\}_{n\in \sigma})$) then for every $\sigma \subset\mathbb{N}$,
	$\{z_n\}_{n\in \sigma}$ is frame for $Z_\sigma.$
	Therefore, for every $\sigma\subset \mathbb{N}$ and for every $k_\sigma \in Z_\sigma$, there exist $0<\alpha_1\leq \beta_1< \infty$ so that, 
	\begin{equation}\label{12}
	\alpha_1 \|k_\sigma\|_{J}^2 \leq \sum\limits_{n \in \sigma} |[k_\sigma, z_n]|^2  \leq \beta_1 \|k_\sigma\|_{J}^2.
	\end{equation}
	Moreover, $\{Pz_n\}_{n \in \sigma^c}$ is a frame for $Z^\perp_{\sigma}.$ Therefore, for every $\sigma\subset \mathbb{N}$ and for every $k^\perp_{\sigma} \in Z^\perp_\sigma$, there exist $0<\alpha_2\leq \beta_2< \infty$ so that
	\begin{equation}\label{13}
		\alpha_2\|k_\sigma^\perp\|_{J}^2 \leq \sum\limits_{n \in \sigma^c} |[k_\sigma^\perp, Pz'_n]|^2  \leq \beta_2 \|k_\sigma^\perp\|_{J}^2.
	\end{equation}
		Furthermore, since $P|Z_{\sigma^c}$ is the orthogonal projection  onto $Z^\perp_{\sigma}$, then from the equation (\ref{13}) we have,
		\begin{equation}\label{14}
		\alpha_2\|k_\sigma^\perp\|_{J}^2 \leq \sum\limits_{n \in \sigma^c} |[k_\sigma^\perp, z'_n]|^2  \leq \beta_2 \|k_\sigma^\perp\|_{J}^2.
	\end{equation}
	Again for every $k \in \mathcal{K}$ we have, $k = k_\sigma + k^\perp_\sigma$.
	Thus applying equations (\ref{12}) and (\ref{14}), for every $\sigma\subset \mathbb{N}$ we obtain,
	
	$$\alpha^*(\|k_\sigma^\perp\|_{J}^2 +\|k_\sigma^\perp\|_{J}^2 )\leq \sum\limits_{n \in \sigma} |[k_\sigma, z'_n]|^2 +\sum\limits_{n \in \sigma^c} |[k_\sigma^\perp, z'_n]|^2  \leq \beta^* (\|k_\sigma^\perp\|_{J}^2 +\|k_\sigma^\perp\|_{J}^2),$$
	where $\alpha^* = \min (\alpha_1, \alpha_2)$	and $\beta^* = \beta_1 + \beta_2$. Consequently, for every $\sigma\subset \mathbb{N}$ and for every $k \in \mathcal{K}$ we obtain,
	
	$$\alpha^*\|k\|_{J}^2  \leq \sum\limits_{n \in \sigma} |[k_\sigma, z'_n]|^2 +\sum\limits_{n \in \sigma^c} |[k_\sigma^\perp, z'_n]|^2  \leq \beta^* \|k\|_{J}^2.$$
This completes the proof.	
\end{proof}

\end{document}
